% Wie schreibe ich das Abstract?
% Bei der Vorbereitung des Abstracts sollten Sie sich an folgenden grundlegenden Punkten orientieren:

%% Motivation des Textes:
%% worin liegt die Bedeutung der entsprechenden Forschung, warum sollte der längere Text gelesen werden?

%% Fragestellung:
%% welche Fragestellung(en) versucht der Text zu beantworten, was ist der Umfang der Forschung, was sind die zentralen Argumente und Behauptungen?

%% Methodologie:
%% welche Methoden/Zugänge nutzt der Autor/die Autorin, auf welche empirische Basis stützt sich der Text?

%% Ergebnisse:
%% zu welchen Ergebnissen kam die Forschung, was sind die zentralen Schlussfolgerungen des Textes?

%% Implikationen:
%% welche Schlussfolgerungen ergeben sich aus dem Text für die Forschung, was fügt der Text unserem Wissen über das Thema hinzu?
 
\begin{abstract}
Diese Arbeit untersucht Verfahren zur kontinuierlichen Gruppierung und zum Clustering von Log-Nachrichten. Solche Verfahren sind ein zentraler Bestandteil von Lösungen zur Überwachung des Systemzustands während der Laufzeit und bei Änderungen, etwa beim Einspielen neuer Software-Releases. Log-Nachrichten weisen häufig eine vergleichbare Struktur auf: Sie bestehen aus festen und variablen Bestandteilen, wobei sich inhaltlich ähnliche Nachrichten im Wesentlichen nur durch unterschiedliche Werte wie IP-Adressen, Zeitstempel oder andere Metadaten unterscheiden.
In einfachen Fällen lassen sich diese variablen Komponenten mithilfe regulärer Ausdrücke extrahieren und die Log-Nachrichten so vereinheitlichen. Diese Methode setzt jedoch eine ständig aktualisierte Sammlung von Regex-Mustern voraus und ist mit erheblichem Wartungsaufwand verbunden. Ziel dieser Arbeit ist daher die Analyse und der Vergleich verschiedener Algorithmen und Methoden zur automatisierten Gruppierung von Log-Nachrichten. Verglichen werden vier algorithmische Verfahren (K-Means, DBSCAN, Drain3, Brain), ein probabilistischer Ansatz auf Basis von MinHash mit Locality-Sensitive Hashing sowie ein vektorbasierter Ansatz, der Sentence-Embeddings in einer PostgreSQL-Datenbank mit der \texttt{pgvector}-Erweiterung persistiert. Die Bewertung erfolgt anhand von drei Kriterien: der Cluster-Qualität, gemessen mit Adjusted Rand Index und V-Measure auf einem gemischten Loghub-2k-Datensatz aus Linux-, Apache- und OpenSSH-Logs ($10\,000$ Nachrichten, $44$ Ground-Truth-Klassen); der Reihenfolge-Stabilität über fünf zufällig permutierte Durchläufe; sowie der strukturellen Eignung für den inkrementellen Betrieb. Auf dem verwendeten Datensatz erzielt der Online-Parser Drain3 mit einem Adjusted Rand Index von $91{,}084\,\%$ und einem V-Measure von $91{,}528\,\%$ die höchsten Qualitätswerte bei zugleich nahezu vollständiger Reihenfolge-Stabilität. Der vektorbasierte Ansatz mit pgvector folgt mit einem Adjusted Rand Index von $79{,}273\,\%$ und einem V-Measure von $88{,}451\,\%$ auf dem zweiten Rang und bleibt insbesondere für Logs mit hoher syntaktischer Variabilität eine sinnvolle Alternative.
\end{abstract}
\cleardoublepage